\chapter{Complexity — Call}
%%% generated by the blueprint pipeline from TCSlib.Complexity.SpaceComplexity.Machines.Call
\section{Overview}

Declarations of the Lean module \texttt{TCSlib.Complexity.SpaceComplexity.Machines.Call}.

\section{Declarations}

\begin{definition}[Clean halting run]
\lean{Complexity.LogProg.CleanRun}
\label{Complexity.LogProg.CleanRun}
\leanok
\uses{Turing.Cfg.init, Turing.MultiTapeTM, Turing.MultiTapeTM.runFrom}
Let $D$ be a deterministic Turing machine with $k$ work tapes, Boolean input and output
symbols, and state set $S$. Given an initial state $q\in S$, a finite Boolean input $V$,
a Boolean $b$, and a natural number $B$, a clean run means that there is a time $T$ at
which the machine has halted, its output is exactly the one-bit list $[b]$, every
work-tape cell is blank, and every work-tape head is at position $0$. At every time
$t\le T$, the absolute position of each work-tape head is at most $B$.
\end{definition}

\begin{lemma}[Concatenating run intervals]
\lean{Complexity.LogProg.runFrom_forall_append}
\label{Complexity.LogProg.runFrom_forall_append}
\leanok
\uses{Turing.Cfg, Turing.MultiTapeTM, Turing.MultiTapeTM.runFrom, Turing.MultiTapeTM.runFrom_add}
\difficulty{4}
Let $M$ be a deterministic multi-tape Turing machine, $c$ a configuration on a fixed
input, $T_1,T_2$ natural numbers, and $Q$ a property of configurations. If $Q$ holds for
every configuration reached from $c$ by at most $T_1$ steps, and also for every
configuration reached from the configuration at time $T_1$ by at most $T_2$ further
steps, then $Q$ holds for every configuration reached from $c$ by at most $T_1+T_2$
steps.
\end{lemma}

\begin{lemma}[Zero offset]
\lean{Complexity.LogProg.off_zero}
\label{Complexity.LogProg.off_zero}
\leanok
\uses{Complexity.LogProg.Seg, Complexity.LogProg.off}
\difficulty{4}
For every finite list of segments, the offset at index zero is zero.
\end{lemma}

\begin{lemma}[Register positions lie in their boxes]
\lean{Complexity.LogProg.regPos_box}
\label{Complexity.LogProg.regPos_box}
\leanok
\uses{Complexity.LogProg.RegBox, Complexity.LogProg.TPos, Complexity.LogProg.TPos.Valid, Complexity.LogProg.callSegs, Complexity.LogProg.regPos, Complexity.LogProg.seg_callSegs_succ, Complexity.LogProg.trackPos_bounds, Complexity.LogProg.wlen}
\difficulty{5}
Let a call specification have an argument list with no repetitions, and let $s$ be less
than the argument-list length. For any track position $tp$ that is valid for segment $s$
of the virtual input, the position assigned to each register by the call-position rule
lies in its register box: registers outside the argument list retain their initial
work-tape positions, while each register in the list has position between $-1$ and the
length of its word, inclusive.
\end{lemma}

\begin{lemma}[Register positions determine the side of a segment]
\lean{Complexity.LogProg.regPos_side}
\label{Complexity.LogProg.regPos_side}
\leanok
\uses{Complexity.LogProg.Consistent, Complexity.LogProg.TPos, Complexity.LogProg.TPos.Valid, Complexity.LogProg.callSegs, Complexity.LogProg.leftSide, Complexity.LogProg.regPos, Complexity.LogProg.seg_callSegs_succ, Complexity.LogProg.trackPos, Complexity.LogProg.wlen}
\difficulty{5}
Let a call specification have an argument list with no repetitions. Let $s$ be a segment
index between $0$ and $m$, let $par$ and $dir$ be Boolean values, and let $tp$ be a
track position consistent with these values for the virtual input segments and valid at
segment $s$. For any $a$ less than both the argument-list length and $m$, if the
position of argument register number $a$ is $-1$, then the side predicate for $s$,
$dir$, and $a$ is true; if its position equals the length of that register's word, then
the side predicate is false.
\end{lemma}

\begin{theorem}[Correctness of a decider call]
\lean{Complexity.LogProg.call_run}
\label{Complexity.LogProg.call_run}
\leanok
\uses{Complexity.LogProg.CSt, Complexity.LogProg.CleanRun, Complexity.LogProg.Consistent, Complexity.LogProg.HaltRel, Complexity.LogProg.RProg, Complexity.LogProg.RegBox, Complexity.LogProg.SimRel, Complexity.LogProg.TPos, Complexity.LogProg.TPos.Valid, Complexity.LogProg.TrackRel, Complexity.LogProg.apply_inputOnly, Complexity.LogProg.apply_moves, Complexity.LogProg.callSegs, Complexity.LogProg.cellOff, Complexity.LogProg.compileTM, Complexity.LogProg.ctr, Complexity.LogProg.inPos, Complexity.LogProg.mkCfg, Complexity.LogProg.mkCfg_state, Complexity.LogProg.off_zero, Complexity.LogProg.regPos, Complexity.LogProg.regPos_box, Complexity.LogProg.regPos_side, Complexity.LogProg.regs_run, Complexity.LogProg.ret2_run, Complexity.LogProg.retPos, Complexity.LogProg.rlen, Complexity.LogProg.runFrom_forall_append, Complexity.LogProg.seam, Complexity.LogProg.seam_eq_mkCfg, Complexity.LogProg.sim_step, Complexity.LogProg.trackPos, Complexity.LogProg.vpos, Complexity.LogProg.vword, Complexity.LogProg.wlen, Complexity.LogProg.wlen_zero_le, Turing.Cfg, Turing.Cfg.init, Turing.FinTM.bufferTape, Turing.FinTM.moveInputPos_neg_val, Turing.MultiTapeTM, Turing.MultiTapeTM.runFrom, Turing.MultiTapeTM.runFrom_add, Turing.MultiTapeTM.runFrom_of_halt, Turing.MultiTapeTM.runFrom_succ_eq_step', Turing.MultiTapeTM.runFrom_zero, Turing.MultiTapeTM.step, Turing.moveInputPos, Turing.moveInputPos_zero}
\difficulty{7}
Let $P$ be a register program with $m$ registers and $d$ available deciders, let $D$ be
a $k_D$-work-tape Boolean Turing machine, and choose its initial state for each decider.
Suppose the initial program configuration is at state $l$, where the call specification
selects decider $j$, a virtual input mode, an ordered list of register arguments, and
the positive- and negative-answer successor states. Let $W$ assign a Boolean word to
each register. Assume each argument register contains its assigned word on a tape
initialized from cell zero, each such tape head is at zero, the argument registers are
pairwise distinct, and the input head is at its initial position. If the selected
decider, started on the virtual input formed from the call specification and the
register words, halts with exactly the single output bit $b$, blank work tapes,
work-tape heads returned to zero, and all work-tape head displacements bounded in
absolute value by $B$ throughout the run, then the compiled machine has a time $T$ at
which it reaches the original program configuration with the input head at its initial
position and the control state set to the successor selected by $b$. At every time
through $T$, each original register head remains in its prescribed region (outside the
argument registers it stays at its initial position; on argument registers it lies
between $-1$ and the word length), and each simulated decider work-tape head remains
within distance $B$ of zero.
\end{theorem}
